\section{Experimental Setup}
\label{sec:setup}

\paragraph{Implementation.}
Experiments run locally on an Apple M4 Pro.  Deep models use PyTorch and
torchvision \cite{paszke2019pytorch} with MPS in FP32; traditional pipelines
use scikit-learn \cite{pedregosa2011scikit} and OpenCV
\cite{bradski2000opencv}.  Data, image, configuration, progress, plotting, and
serialisation utilities use NumPy, pandas, tqdm, scikit-image, Pillow,
Matplotlib, joblib, torchvision, and PyYAML
\cite{harris2020array,mckinney2010data,costaluis2019tqdm,vanderwalt2014scikitimage,
clark2026pillow,hunter2007matplotlib,joblib2025,marcel2010torchvision,pyyaml2025}.
Random seeds for Python, NumPy, data samplers,
workers, and model initialisation are fixed to 42.  Data decoding applies EXIF
orientation and handles grayscale or alpha-channel inputs by safe RGB
conversion.  Failed or unreadable images are fatal rather than silently
skipped.

Scratch ResNet18 uses a 25-epoch cap and patience 6.  The two pretrained
ResNet18 augmentation variants share initial pretrained weights, head
initialisation, sampler order, optimiser, scheduler, 15-epoch cap, and
patience 4; only the training transform changes.  EfficientNet-B0 uses the
same strong-augmentation optimisation policy.  A disposable preflight chooses
a feasible micro-batch without contributing weights or measurements to the
canonical run.  Before epoch one, the model is rebuilt and all random sources
are reseeded.

\paragraph{Metrics.}
For each of 500 labels we compute precision, recall, and F1 with zero assigned
when a denominator is empty, then report their unweighted macro averages.
Top-$k$ accuracy counts an example as correct when its true class is among the
$k$ largest scores.  Score columns are aligned through the estimator's
explicit class order, and stable descending ranking breaks exact ties in favour
of the lower label.  Numeric confusion matrices use rows for true labels and
columns for predictions.  We retain both overall/top-1 accuracy names for
rubric compatibility, but they are the same statistic.

Traditional training time surrounds classifier \texttt{fit} only, and its
test time surrounds scoring an already-loaded cached feature matrix.  Deep
training time includes loading, augmentation, transfer, forward/backward, and
optimiser work in the training loop, excluding validation and serialisation.
Deep test time is synchronised model-forward time; full evaluation wall time
is also logged.  Consequently, traditional and deep timing numbers describe
different pipeline boundaries and should not be interpreted as a direct
hardware-speed contest.

\paragraph{Selection and reproducibility.}
For Stage~3, validation metrics are independently recomputed from each best
checkpoint before freeze.  Test selection includes scratch ResNet18, the
validation-best pretrained ResNet18 augmentation, and completed
EfficientNet-B0; the basic-augmentation and frozen-head ResNet18 runs remain
validation-only.  For rewritten Stage~2, the selection manifest binds the 19
validation configurations, selected component order, normalisation, model and
cache hashes, train-fitted BoVW vocabulary, and Stage~1 split hashes.  Test
evaluation rejects a missing or mismatched manifest, an unlisted model, and
silent replacement of canonical metrics.  This freezes the current state
going forward.  Because no independently verifiable pre-test manifest survives
for the already-recorded traditional results, it does not establish historical
chronology.  Atomic caches, checkpoints, predictions, metric files, and
configuration hashes nevertheless make the evidence currently reported
traceable to its bound source.
Figure~\ref{fig:convergence} shows observed optimisation curves reconstructed
from the persisted epoch logs, without rerunning inference.

\begin{figure*}[t]
  \centering
  \includegraphics[width=0.92\textwidth]{training_convergence.png}
  \caption{Training loss and validation macro F1 for scratch ResNet18 and the
  best pretrained EfficientNet-B0.  Transfer learning converges faster and to
  a substantially stronger validation solution.}
  \label{fig:convergence}
\end{figure*}
